\documentclass[../../../../main.tex]{subfiles}
\begin{document}

\begin{flushright}
\footnotesize\textit{【自動文字起こし・要確認】}
\end{flushright}

[解] $x=0, \pi, \pi/2$での成立が必要。又、$0 \le C < 2\pi$とする.
\begin{align}
  3a + b(1+2\cos C - 3\sin C) &= 1 \label{eq2_1} \\
  -a + b(1-2\cos C + 3\sin C) &= 1 \label{eq2_2} \\
  4a + b(1+2\sin C + 3\cos C) &= 1 \label{eq2_3}
\end{align}
以下$\sin C = \alpha$, $\cos C = \beta$とおく. \eqref{eq2_1}+\eqref{eq2_2}から
\begin{equation}
  a + b = 1 \label{eq2_4}
\end{equation}
\eqref{eq2_4}を\eqref{eq2_1}, \eqref{eq2_3}に代入
\begin{align}
  b(-2 + 2\beta - 3\alpha) + 2 &= 0 \label{eq2_5} \\
  b(-3 + 3\beta + 2\alpha) + 3 &= 0 \label{eq2_6}
\end{align}
\eqref{eq2_5}$\times 3$ - \eqref{eq2_6}$\times 2$から
\[
  b[-6 + 6\beta - 9\alpha + 6 - 6\beta - 4\alpha] = 0
\]
\[
  \alpha b = 0
\]
したがって, $b=0$又は$\alpha=0 \iff C=0, \pi$ ($\because 0 \le C < 2\pi$) である.

$1^\circ \ b=0$の時\\
$af(x)=1$が任意の$x$で成立するような$a$は存在せず, 不適.

$2^\circ \ C=0$の時 (この時$\alpha=0, \beta=1$)\\
\eqref{eq2_5}から$2=0$となり不適

$3^\circ \ C=\pi$の時 ($\alpha=0, \beta=-1$)\\
\eqref{eq2_5}から$b=\frac{1}{2}$, \eqref{eq2_4}から$a=\frac{1}{2}$である. 逆にこの時
\[
  \frac{1}{2}f(x) + \frac{1}{2}f(x-\pi) = 1
\]
は任意の$x$について成立し, 十分

以上を, $C$を一般角に直して
\[
  (a, b, C) = \left(\frac{1}{2}, \frac{1}{2}, \pi + 2n\pi\right) \ (n \in \mathbb{Z})
\]

\end{document}
